%=============================================================================!
% 5.3  THE ANGULAR MOMENTUM OF MANY BODIES                                    !
%=============================================================================!
\section{The angular momentum of many bodies}
\label{sec:ro-system}

\Cref{sec:ne-bodies} added the momenta of a group of bodies and found that
the forces between them cancel in pairs, so that only outside forces change
the total. This section does the same for angular momentum, and finds one
extra condition: the forces between two bodies must act along the line
joining them.

\subsection*{Internal torques cancel}

For $N$ bodies, at $\vb{r}_i$ with momenta $\vb{p}_i$, the total angular
momentum about the origin is
\begin{equation*}
   \vb{L} = \sum_{i=1}^{N}\vb{r}_i\times\vb{p}_i ,
\end{equation*}
and by \cref{eq:ro-torque} for each body, $\dd\vb{L}/\dd t = \sum_i\vb{r}_i
\times\vb{F}_i$, with $\vb{F}_i$ the total force on body $i$. As in
\cref{sec:ne-bodies}, split each force into the external force
$\vb{F}_i^{\mathrm{ext}}$, from outside the group, and the forces
$\vb{F}_{ij}$ on body $i$ from the other bodies $j$. The torques of the forces between bodies $i$ and $j$ appear as a
pair,
\begin{align*}
   \vb{r}_i\times\vb{F}_{ij} + \vb{r}_j\times\vb{F}_{ji}
   &= \vb{r}_i\times\vb{F}_{ij} - \vb{r}_j\times\vb{F}_{ij}
   && \text{by the third law, $\vb{F}_{ji} = -\vb{F}_{ij}$,}\\
   &= (\vb{r}_i - \vb{r}_j)\times\vb{F}_{ij}
   && \text{a sum splits, in the first vector as in the second.}
\end{align*}
The third law alone does not make this zero: two equal and opposite
forces side by side, like the pushes on the book, form a couple. But if
the force between two bodies acts along the line joining them,
$\vb{F}_{ij} = c\,(\vb{r}_i - \vb{r}_j)$ for some number $c$, the pair
gives $c\,(\vb{r}_i - \vb{r}_j)\times(\vb{r}_i - \vb{r}_j) = \vb{0}$. Then
every pair drops out, and
\begin{equation}
   \frac{\dd\vb{L}}{\dd t} = \sum_{i=1}^{N}\vb{r}_i\times
   \vb{F}_i^{\mathrm{ext}} = \vb{G}^{\mathrm{ext}} :
   \label{eq:ro-system}
\end{equation}
only the external torque changes the total angular momentum, and when it is
zero, the angular momentum of the group is conserved.

\subsection*{Forces between atoms}

Atoms whose potential energy is a sum of terms $\varphi(r_{ij})$, each
depending only on the distance $r_{ij}$ between two atoms, push and pull
along the lines joining them. Differentiating $r_{ij}^2 = (x_j - x_i)^2 +
(y_j - y_i)^2 + (z_j - z_i)^2$ with respect to $x_i$ by the chain rule
gives $2r_{ij}\,\partial r_{ij}/\partial x_i = -2(x_j - x_i)$, so
$\partial r_{ij}/\partial x_i = -(x_j - x_i)/r_{ij}$, and likewise for $y_i$
and $z_i$. The force on atom $i$ from this term is then
\begin{equation*}
   -\grad_i\varphi(r_{ij}) = -\varphi'(r_{ij})\,\grad_i r_{ij}
   = \varphi'(r_{ij})\,\frac{\vb{r}_j - \vb{r}_i}{r_{ij}} ,
\end{equation*}
along the line from atom $i$ to atom $j$, as \cref{ex:en-pair} found. An
isolated group of such atoms therefore conserves its angular momentum as
well as its momentum. Many of the forces of Chapter~8 depend on three or more atoms
at once and cannot be split into pairs, but Chapter~7 shows that, whenever
the potential energy is unchanged when the whole group is turned, the
internal torques still add to zero.

\subsection*{Turning about the centre of mass}

Unlike momentum, angular momentum depends on the point it is measured
about: about the point $\vb{a}$, each position is measured from $\vb{a}$,
so that $\sum_i(\vb{r}_i - \vb{a})\times\vb{p}_i$ takes the place of
$\vb{L}$. Measure each position and velocity from the centre of mass
$\vb{R}$ and its velocity $\vb{V} = \vb{P}/M$ (\cref{sec:ne-bodies}),
$\vb{r}_i = \vb{R} + \vb{r}_i'$ and $\vb{v}_i = \vb{V} + \vb{v}_i'$. The
primed quantities satisfy $\sum_i m_i\vb{r}_i' = \sum_i m_i\vb{r}_i -
M\vb{R} = \vb{0}$ and, differentiating, $\sum_i m_i\vb{v}_i' = \vb{0}$.
Then
\begin{align*}
   \vb{L} &= \sum_i m_i\,(\vb{R} + \vb{r}_i')\times(\vb{V} + \vb{v}_i')\\
   &= \Bigl(\sum_i m_i\Bigr)\vb{R}\times\vb{V}
      + \vb{R}\times\sum_i m_i\vb{v}_i'
      && \text{multiplying out}\\
   &\quad + \Bigl(\sum_i m_i\vb{r}_i'\Bigr)\times\vb{V}
      + \sum_i m_i\,\vb{r}_i'\times\vb{v}_i'\\
   &= \vb{R}\times\vb{P} + \vb{L}' ,
   \qquad \vb{L}' = \sum_i m_i\,\vb{r}_i'\times\vb{v}_i'
   && \text{$M\vb{V} = \vb{P}$; the middle two vanish.}
\end{align*}
The angular momentum of a group is that of its total momentum carried at
the centre of mass, plus $\vb{L}'$, its angular momentum about its own
centre of mass, which does not depend on where the origin is, since moving the origin
shifts each $\vb{r}_i$ and $\vb{R}$ alike and leaves $\vb{r}_i' = \vb{r}_i
- \vb{R}$ unchanged. For a
thrown spinning ball, the first part belongs to the flight and the second
to the spin; for a molecule, $\vb{L}'$ is the turning of the molecule
itself. For an isolated group both parts are conserved separately: $\vb{P}$
is constant, so $\dd(\vb{R}\times\vb{P})/\dd t = \vb{V}\times\vb{P} =
\vb{0}$, since $\vb{P} = M\vb{V}$ is parallel to $\vb{V}$; and $\vb{L}$ is
constant by \cref{eq:ro-system}, so $\vb{L}' = \vb{L} - \vb{R}\times\vb{P}$
is constant too.

\begin{keyidea}
Forces between pairs of bodies along the lines joining them cancel in the
total torque, so only external torques change the angular momentum of a
group, and an isolated group of atoms keeps it. The total splits into the
part carried by the centre of mass and the part $\vb{L}'$ about it.
\end{keyidea}

\begin{selfcheck}
Two bodies of \SI{1}{\kilogram} at $(\pm0.5, 0, 0)$\,\si{\metre} move with
velocities $(0, \mp1, 0)$\,\si{\metre\per\second}. Find their total
momentum and their angular momentum about the origin and about the point
$(0, 1, 0)$\,\si{\metre}.
\end{selfcheck}
